\chapter{Implementation}
\label{chap:implementation}

This work is an add-on contribution to the Darwin G\"odel Machine paper,
whose code is public.\footnote{\texttt{https://github.com/jennyzzt/dgm.git}}
We forked the repo and added an advanced scheduler mechanism to try to
get as much as possible out of this system. Each scheduler implementation
is in the \texttt{schedulers.py} file. This chapter documents the fork
additions (Section~\ref{sec:fork}), the scheduler implementations
(Section~\ref{sec:scheduler-impl}), and the operational guards that keep
long runs robust (Section~\ref{sec:guards}).

Figure~\ref{fig:components} gives a high-level component view of the
forked system and how the parts interact across a single generation.
The evolution loop in \texttt{DGM\_outer.py} drives a pluggable
scheduler, which delegates child generation and evaluation to
\texttt{self\_improve\_step.py}. Generation flows through the coding
agent and the LLM stack out to the OpenRouter API; evaluation flows
through the benchmark registry and the Docker utilities into a
per-task SWE-bench harness container. Results from both flows are
recorded per child and then appended to \texttt{dgm\_metadata.jsonl},
which feeds the archive consulted by the scheduler in the next
generation.

\begin{figure}[tb]
  \centering
  \resizebox{\textwidth}{!}{%
  \begin{tikzpicture}[
      font=\small\ttfamily,
      every node/.style={align=center},
      mod/.style={draw, rounded corners=2pt, fill=black!3,
                  minimum width=32mm, minimum height=9mm},
      new/.style={draw, rounded corners=2pt, fill=blue!8,
                  line width=0.8pt,
                  minimum width=32mm, minimum height=9mm},
      ext/.style={draw, dashed, rounded corners=2pt, fill=black!2,
                  minimum width=32mm, minimum height=9mm},
      data/.style={draw, rounded corners=2pt, fill=orange!12,
                   minimum width=32mm, minimum height=9mm},
      flow/.style={-{Latex[length=2mm]}, thick},
      gen/.style={flow, draw=blue!60!black},
      eval/.style={flow, draw=red!50!black},
      lbl/.style={font=\scriptsize\rmfamily, fill=white, inner sep=1pt},
  ]

  % Backbone (centered)
  \node[mod]  (outer)    at ( 0,    0)    {DGM\_outer.py\\\scriptsize evolution loop};
  \node[new]  (sched)    at ( 0,   -1.7)  {schedulers.py\\\scriptsize Base / Baseline / Hyperband / ASHA / GA};
  \node[mod]  (sis)      at ( 0,   -3.4)  {self\_improve\_step.py\\\scriptsize child generation + evaluation};

  % Left column (generation pipeline)
  \node[mod]  (agent)    at (-3.2, -5.4)  {coding\_agent.py\\\scriptsize patch generation};
  \node[new]  (gaprompt) at (-7.8, -5.4)  {prompts/\\ga\_mutation\_prompt.py};
  \node[mod]  (llmwt)    at (-3.2, -7.1)  {llm\_withtools.py\\\scriptsize tool-calling layer};
  \node[mod]  (llm)      at (-3.2, -8.8)  {llm.py\\\scriptsize base LLM caller};
  \node[ext]  (or)       at (-3.2,-10.5)  {OpenRouter API\\\scriptsize minimax-m2.5};

  % Right column (evaluation pipeline)
  \node[new]  (bench)    at ( 3.2, -5.4)  {benchmarks/config.py\\\scriptsize task pool registry};
  \node[mod]  (docker)   at ( 3.2, -7.1)  {utils/docker\_utils.py\\\scriptsize container provisioning};
  \node[ext]  (dock)     at ( 3.2, -8.8)  {Docker daemon};
  \node[ext]  (swe)      at ( 3.2,-10.5)  {SWE-bench harness\\\scriptsize per-task container};

  % Data artifacts (centered, below both columns)
  \node[data] (meta)     at ( 0,   -12.4) {dgm\_metadata.jsonl\\\scriptsize per-generation log};
  \node[data] (archive)  at ( 0,   -14.0) {output\_dgm/\\\scriptsize archive of compiled children};

  % Control flow (backbone)
  \draw[flow] (outer) -- node[lbl, right]{select scheduler} (sched);
  \draw[flow] (sched) -- node[lbl, right]{run\_generation()} (sis);

  % Generation branch
  \draw[gen] (sis.south) -- ++(0,-0.35) -| (agent.north);
  \draw[gen] (agent) -- (llmwt);
  \draw[gen] (llmwt) -- (llm);
  \draw[gen] (llm) -- (or);
  \draw[gen] (gaprompt.east) -- node[lbl, above, yshift=5mm]{GA only} (agent.west);

  % Evaluation branch
  \draw[eval] (sis.south) -- ++(0,-0.35) -| (bench.north);
  \draw[eval] (bench) -- (docker);
  \draw[eval] (docker) -- (dock);
  \draw[eval] (dock) -- (swe);

  % Branch terminals -> metadata (meta sits between or and swe, no archive crossings)
  \draw[flow] (or.south)  |- (meta.west);
  \draw[flow] (swe.south) |- (meta.east);
  \draw[flow] (meta) -- (archive);

  % Feedback arrow routed via far-left margin, well clear of all nodes
  \draw[flow, dashed]
    (archive.west) -- ++(-9.5,0) coordinate (fbk)
    -- (fbk |- sched.west)
    -- (sched.west);
  \node[lbl] at ($(fbk |- sched.west) + (4mm, -3mm)$) {next generation};

  % Legend (top-left, above outer)
  \node[new,  minimum width=22mm, minimum height=5mm, font=\scriptsize\ttfamily] (lg1) at (-7.8, 1.7) {new file};
  \node[mod,  minimum width=22mm, minimum height=5mm, font=\scriptsize\ttfamily] (lg2) at (-5.1, 1.7) {modified file};
  \node[ext,  minimum width=22mm, minimum height=5mm, font=\scriptsize\ttfamily] (lg3) at (-2.4, 1.7) {external};
  \node[data, minimum width=22mm, minimum height=5mm, font=\scriptsize\ttfamily] (lg4) at ( 0.3, 1.7) {data artifact};

  \end{tikzpicture}%
  }
  \caption{Component view of the forked DGM system. Blue-filled boxes
    are files added by this work, grey-filled boxes are files of the
    upstream DGM that were modified, dashed boxes are external
    components (Docker daemon, OpenRouter, SWE-bench harness), and
    orange boxes are persistent data artifacts. Blue arrows trace the
    child-generation pipeline and red arrows trace the evaluation
    pipeline; both converge on \texttt{dgm\_metadata.jsonl}, which
    feeds the archive consulted by the scheduler in the next
    generation.}
  \label{fig:components}
\end{figure}

\section{Fork Additions}
\label{sec:fork}

We forked the Darwin G\"odel Machine repo and modified several files from
the original repo:

\begin{itemize}
  \item \textbf{\texttt{DGM\_outer.py}}: changed the default benchmark to
    SWE-bench Verified Mini; created a scheduler integration with a
    pluggable scheduler hook, re-evaluation-on-promotion budget, new
    metrics, and the ASHA/GA wiring; added the
    \texttt{--generation\_task\_budget\_total} CLI argument
    (\texttt{DGM\_outer.py:277}, default \texttt{None}; the value
    \texttt{100} used in the experiments is a run-time flag, not a code
    default).
  \item \textbf{\texttt{self\_improve\_step.py}}: four big changes. Added
    support for the OpenRouter API key and new models; added the new
    benchmark switch, scheduler budget/metrics, and the GA mutation hook;
    wired the argparse guard. (Note: the Docker repo-copy and container
    git-config fixes are \emph{called} from here but live in
    \texttt{utils/docker\_utils.py}, see Section~\ref{sec:guards}.).
  \item \textbf{\texttt{coding\_agent.py}}: scheduler budget and metrics;
    the empty-patch retry loop (Section~\ref{sec:guards}).
  \item \textbf{\texttt{llm.py}}: the base LLM caller layer. Added the
    model registry entries \texttt{openrouter/minimax/minimax-m2.5} and
    \texttt{openrouter/openai/gpt-5-mini} (\texttt{llm.py:20-21}) as
    usable models, and adapted the layer to keep everything in sync with
    the original code. Minimax is a thinking model, so part of its output
    is a thinking segment wrapped in thinking/JSON tags that must be
    stripped to prevent empty-response failures.
  \item \textbf{\texttt{llm\_withtools.py}}: the tool-calling layer.
    Replaced the default model GPT-5-Mini with Minimax. Created a
    dedicated \texttt{chat.completion} path for Minimax with its standard
    \texttt{tool\_calls} format and no strict mode, because the OpenRouter
    \texttt{chat.completion} schema does not support strict tool-calling.
    It now preserves the full tool-call objects in history, catches
    malformed JSON to skip gracefully without crashing, and mirrors the
    tag-stripping helpers from \texttt{llm.py}.
\end{itemize}

We created new files for the new scheduler mechanisms:

\begin{itemize}
  \item \textbf{\texttt{schedulers.py}}: the entire pluggable scheduler
    layer, containing \texttt{BaseScheduler} plus
    Baseline/Hyperband/ASHA/GA subclasses, per-generation budget
    recording, and timeout constants. This file did not exist in the
    original DGM repo; DGM only had the fixed three-stage evaluation that
    the Baseline scheduler reproduces. This is the core thesis
    contribution: all scheduler logic resides here and the schedulers are
    swappable from one file.
  \item \textbf{\texttt{benchmarks/config.py}}: a configurable test suite
    to select SWE-bench, Polyglot, or SWE-bench Verified Mini. Replaces
    DGM's hardcoded SWE-bench paths. Created because we needed
    reproducible pinned SWE-bench Verified Mini subsets and cumulative
    rung sizes for scheduler evaluation.
  \item \textbf{\texttt{prompts/ga\_mutation\_prompt.py}}: the mutation
    prompt for the genetic algorithm, with no error-log context, asking
    the LLM for random exploration mutations. Reuses the coding-agent
    summary and current-code logic from the self-improvement prompts and
    applies high temperature at call time.
  \item \textbf{\texttt{compare\_scheduler\_runs.py}}: reads all results
    from the \texttt{dgm\_metadata.jsonl} files of the four scheduler runs
    and aggregates them into one JSON for cross-scheduler comparison.
  \item \textbf{\texttt{tests/test\_benchmark\_registry.py}}: verifies the
    pinned, cumulative benchmark subset and keeps the legacy benchmarks
    under a legacy register key for backward compatibility. Added because
    a fixed subset size is required for a fair comparison between
    schedulers.
\end{itemize}

\section{Scheduler Implementations}
\label{sec:scheduler-impl}

\subsection{BaseScheduler}

This is the parent scheduler all others override. It defines the contract
abstractly: \texttt{run\_generation(...)} and
\texttt{get\_generation\_child\_count(...)} both
\texttt{raise NotImplementedError}, so subclasses must supply the
implementation. It provides two shared parallel-executor helpers:
\texttt{\_generate\_children(...)} for parallel child-patch generation and
\texttt{\_evaluate\_children(...)} for parallel evaluation. Promotion,
ranking, and metadata-marking are shared module-level helpers
(\texttt{\_promote\_top\_k}, \texttt{\_rank\_children},
\texttt{\_mark\_promotion\_decisions}) that every subclass calls.

The \texttt{BaseScheduler} contract and shared helpers are listed in
Listing~\ref{lst:base-scheduler}
(Appendix~\ref{app:scheduler-pseudocode}).

\subsection{Baseline}

Builds on BaseScheduler and adds five instance attributes: three task-ID
lists \texttt{stage\_small\_issues} / \texttt{stage\_medium\_issues} /
\texttt{final\_stage\_issues} (one per eval stage),
\texttt{expected\_task\_counts} (used to check whether a child ran enough
tasks for its results to be valid), and \texttt{test\_more\_threshold}
(the promotion gate; a child advances if it solves at least this fraction
of the stage). \texttt{get\_generation\_child\_count} returns how many
children to spawn per generation. \texttt{run\_generation} is a linear,
no-loop implementation with three sequential eval blocks; a
\texttt{consumed\_budget} accumulator increments after each block by
number of children $\times$ stage size.

The stage structure is: \textbf{Stage 0 (small)} evaluates all children,
promoting those that resolve at least 40\% of small-stage tasks;
\textbf{Stage 1 (medium)} takes the Stage-0 survivors and promotes those
whose score is greater than or equal to the full-evaluation threshold;
\textbf{Stage 2 (full)} takes the remaining children, evaluates them on the
full suite, and keeps all of them. Every stage promotion has two steps:
(a) a quality gate filtered by the threshold (Stages 0 and 1), then (b) a
budget cap that takes the max number of promotions equal to remaining
budget divided by next-stage size, keeping the top-k ranked. The gate
selects eligible candidates; the budget cap then trims. All scheduler
metadata is saved via \texttt{finalize\_child\_metadata} per child.

The full \texttt{run\_generation} pseudocode for Baseline is given in
Listing~\ref{lst:baseline-pseudo}
(Appendix~\ref{app:scheduler-pseudocode}).

\subsection{Hyperband}

Builds on BaseScheduler and adds three instance attributes:
\texttt{self.eta} (reduction factor; each rung keeps only the top
$1/\eta$ and kills the rest), \texttt{self.budgets} (ascending list of
task counts per rung), and \texttt{self.reference\_instance\_ids} (fixed
task-ID pool, for a consistent pool across children). Three helpers:
\texttt{\_generation\_budget\_target()} defines the target task count per
generation; \texttt{\_budget\_cost(n)} simulates the full SHA cost under
an incremental task-cost model;
\texttt{get\_generation\_child\_count} returns
\texttt{--hyperband\_initial\_children} directly when that flag is set,
and only when it is unset falls back to picking the largest \texttt{n}
with \texttt{\_budget\_cost(n) <= target}. The experiments used
\texttt{--hyperband\_initial\_children 15}, so the
\texttt{\_budget\_cost} search path was not exercised in the actual runs.

\texttt{\_task\_increments(generation\_seed)} does two things per
generation: it shuffles the task list with a seed so all children get the
same shuffle in a generation, then cuts the shuffled list into chunks, one
per rung. One implementation detail worth noting: each chunk holds
\emph{different} tasks, so a child at rung~1 is tested only on tasks
3 to 10, not tasks 1 to 10. This is a slight deviation from textbook
Hyperband.

\texttt{run\_generation} is implemented for the SWE-bench benchmark only.
The loop is over rungs, each with a \texttt{budget\_size} and a
\texttt{task\_increment}. The synchronous barrier means
\texttt{\_evaluate\_children} returns only when all current candidates are
done, so every candidate finishes a rung before any promotion. Consumed
budget is the sum of task increments $\times$ candidates, so it counts
delta tasks only. Promotion works as follows: at the last rung all
survivors pass; otherwise promote $\lceil$candidates $/ \eta\rceil$ of the
top-K with score greater than 0. All rung data is stored in the per-rung
dict, and \texttt{current\_candidates = promoted} feeds the next rung.

The main differences from Baseline: the winner is the single top finalist
(\texttt{winner = \_rank\_children(survivors)[0]}); in
\texttt{finalize\_child\_metadata} only the winner gets
\texttt{post\_improve\_diagnose = args.post\_improve\_diagnose} (others
\texttt{False}), saving cost so the diagnosis runs once; and the archive
is updated with a single member, so only the winner is added each
generation. Five structural contrasts vs Baseline: (1) Baseline archives
all compiled, fully-evaluated children while Hyperband archives only the
winner, so archive growth rate differs; (2) Hyperband's rung structure is
configurable while Baseline's is three hard-coded blocks; (3) Hyperband
promotes with \texttt{require\_positive=True} while Baseline uses a static
\texttt{resolved >= 0.4 * small} gate, a different criterion shape; (4)
post-improvement diagnosis is winner-only in Hyperband, all-children in
Baseline; (5) the ``fully tested'' task count is read from
\texttt{budgets[-1]} in Hyperband versus
\texttt{expected\_task\_counts[-1]} in Baseline.

The full \texttt{run\_generation} pseudocode for Hyperband is given in
Listing~\ref{lst:hyperband-pseudo}
(Appendix~\ref{app:scheduler-pseudocode}).

\subsection{ASHA}

ASHA has the same subclass attributes as Hyperband (\texttt{eta} /
\texttt{budgets} / \texttt{reference\_instance\_ids}) and byte-identical
\texttt{\_generation\_budget\_target} / \texttt{\_budget\_cost} /
\texttt{get\_generation\_child\_count} / \texttt{\_task\_increments}
helpers.

The unique points: it precomputes a quota array at startup,
\texttt{promote\_quotas[r] = max(1, ceil(n\_children / eta\^{}(r+1)))},
before the loop, differing from Hyperband which recomputes the promote
count live per rung. ASHA fixes quotas upfront off the initial
\texttt{n\_children}, and the quota is a \emph{cumulative} cap on total
promotions ever from rung~r, not a per-rung count. There is no
synchronous barrier: a single shared \texttt{ThreadPoolExecutor} serves
the whole generation, all children are submitted at rung~0, and
\texttt{wait(FIRST\_COMPLETED)} processes evals as each lands. This
removes one of Hyperband's barriers. ASHA does not wait for all children
to finish a rung; it promotes each child as it finishes a rung if it
meets the promotion threshold, without waiting for the others.

ASHA's promote-on-arrival structure still ranks candidates by score, but
because the quota fills from the completions seen so far, a fast mediocre
child can take a quota slot before a slower but better child has landed, a
trade-off of ASHA. It refills slots immediately: each promotion
synchronously submits the candidate's next-rung evaluation within
\texttt{\_asha\_try\_promote}, folding the new future into the pending set
so workers never idle between rungs. Finally, ASHA uses a longer timeout
(\texttt{ASHA\_GENERATION\_EVAL\_TIMEOUT\_SECONDS}, 36~h, versus the
1.5~h patch-generation timeout) because its rung chains run for longer
and need more time to terminate the full evaluation; this separation of
the eval-loop timeout from the patch-generation timeout was a dedicated
fix.

The full \texttt{run\_generation} pseudocode for ASHA is given in
Listing~\ref{lst:asha-pseudo}
(Appendix~\ref{app:scheduler-pseudocode}).

\subsection{GA}

The shortest scheduler. It inherits all of BaselineScheduler and adds one
attribute, \texttt{self.mutation\_temperature =
args.ga\_mutation\_temperature}. It overrides one method,
\texttt{\_generate\_children}: it ignores caller flags and calls
\texttt{super().\_generate\_children(..., blind\_mutation=True,
mutation\_temperature=self.mutation\_temperature)}. The mutation logic
itself lives in \texttt{self\_improve\_step.py:470} and
\texttt{prompts/ga\_mutation\_prompt.py}; error-log diagnosis is replaced
by blind high-temperature sampling. Selection is tournament parent-choice
in \texttt{DGM\_outer.py}, not in the scheduler.

The GA subclass and its single override are given in
Listing~\ref{lst:ga-pseudo}
(Appendix~\ref{app:scheduler-pseudocode}).

\section{Operational Guards}
\label{sec:guards}

Long DGM runs are sensitive to a class of silent failures in which a
child agent fails to evaluate but still consumes its full budget. We add
three guards to keep these from polluting the comparison.

\paragraph{Argparse guard.}
After a child patch is applied inside its container, the harness runs
\texttt{python /dgm/coding\_agent.py --help} under a 30-second timeout
before any evaluation begins. The help output is checked for every CLI
flag the evaluation drivers pass: \texttt{--problem\_statement},
\texttt{--git\_dir}, \texttt{--chat\_history\_file},
\texttt{--base\_commit}, \texttt{--outdir}, \texttt{--test\_description},
\texttt{--self\_improve}, plus \texttt{--instance\_id} for SWE-style
benchmarks or \texttt{--language} for polyglot. If the help invocation
does not exit, or any required flag is absent, the child is marked
\texttt{model\_patch\_notempty = False} with a
\texttt{compile\_failure\_reason}, its \texttt{budget\_history} is
emptied, and it is excluded from evaluation. This prevents a
self-modification that removes or renames a required argument from
silently consuming the full task budget on argparse-exit-2 failures. The
same constraint is also stated explicitly in the self-improvement and
GA-mutation prompts, reducing how often a mutation strips the flags in
the first place.

\paragraph{Empty-patch retry.}
The coding agent's \texttt{forward()} method wraps generation in a
validation-retry loop. After the agent runs, \texttt{validate\_patch()}
rejects a diff that is empty/whitespace-only or that touches only files
under \texttt{test/} or \texttt{tests/}. On rejection the working tree is
reset to the base commit, the instruction is rebuilt with an explicit
failure reason and stronger wording (``your previous attempt FAILED
because\ldots''), and the agent is re-invoked with a fresh message
history; this repeats up to \texttt{max\_retries} (default 3). Token
usage is accumulated across attempts. If all attempts still produce an
invalid patch the failure is logged and generation proceeds, with the
empty patch handled downstream. The guard targets models that describe
edits in prose instead of applying them via the editor/bash tools
(observed with MiniMax). Important to have in mind that retries do not
vary sampling temperature, the corrective signal is the strengthened
prompt and context reset, not increased randomness.

\enlargethispage{3\baselineskip}
\paragraph{Docker startup-retry and per-image build lock.}
Container provisioning in \texttt{build\_dgm\_container()} has two guards.
First, image builds are serialized per image name: a module-level dict of
\texttt{threading.Lock} objects, itself protected by a global guard lock,
returns one lock per \texttt{image\_name}. The build runs while that lock
is held, and image existence is re-checked inside the lock (double-checked
locking), so under concurrent child generation only one thread builds a
given image while the rest skip straight to reuse. A build exception
returns \texttt{None}. Second, container startup is retried with
exponential backoff: for up to \texttt{max\_retries} attempts (default 3)
the code starts the container, syncs the repository, and resets its git
state; on failure it logs the error, removes any half-started container so
the name is freed, waits $10 \cdot 2^{(\text{attempt}-1)}$ seconds (10~s,
20~s, 40~s), and retries. After exhausting attempts it logs the last
error and returns \texttt{None}. The retry exists because the Docker
daemon transiently fails under concurrent container-startup pressure, this
was observed on consumer hardware with five or more parallel evaluation
containers contending for the same daemon.
